% ============================================================
% APUNTES DE ESTUDIO — DERECHO PROCESAL CIVIL Y COMERCIAL
% La audiencia preliminar, la prescripción adquisitiva
% y los oficios judiciales
% Documento generado a partir de apuntes de clase (18/09/2026)
% ============================================================

\documentclass[12pt,oneside]{book}

% ------------------------------------------------------------
% METADATOS DEL DOCUMENTO
% ------------------------------------------------------------
\newcommand{\universidad}{Universidad de Buenos Aires (UBA)}
\newcommand{\materia}{Derecho Procesal Civil y Comercial}
\newcommand{\titulo}{La audiencia preliminar, la prescripción adquisitiva y los oficios judiciales}
\newcommand{\autor}{Isaac Edgar Camacho Ocampo}
\newcommand{\anio}{2026}
\newcommand{\reflexionmotivacional}{Comprender el recorrido de un expediente real es el primer paso para ejercer con criterio.}

% ------------------------------------------------------------
% PAQUETES
% ------------------------------------------------------------
\usepackage[utf8]{inputenc}
\usepackage[T1]{fontenc}
\usepackage[spanish,es-nodecimaldot]{babel}

\usepackage[
    top=2.5cm,
    bottom=2.5cm,
    left=3cm,
    right=2.5cm,
    headheight=15pt
]{geometry}

\usepackage{amsmath}
\usepackage{amssymb}
\usepackage{mathtools}

\usepackage{graphicx}
\usepackage{float}
\usepackage{caption}
\usepackage{subcaption}

\usepackage{booktabs}
\usepackage{array}
\usepackage{longtable}
\usepackage{tabularx}

\usepackage{enumitem}

\usepackage{xcolor}
\usepackage[most]{tcolorbox}

\usepackage{fancyhdr}
\usepackage{ragged2e}

\usepackage[
    colorlinks=true,
    linkcolor=blue!50!black,
    citecolor=green!40!black,
    urlcolor=blue!60!black,
    bookmarks=true
]{hyperref}

\hypersetup{
    pdftitle={\titulo},
    pdfauthor={\autor},
    pdfsubject={Apuntes universitarios}
}

% ------------------------------------------------------------
% RUTAS DE IMÁGENES
% ------------------------------------------------------------
\graphicspath{
    {./img/}
    {./graficos/}
    {./figuras/}
}

% ------------------------------------------------------------
% CONFIGURACIÓN DE ÍNDICE, PÁRRAFOS Y LISTAS
% ------------------------------------------------------------
\setcounter{tocdepth}{2}

\setlength{\parindent}{0pt}
\setlength{\parskip}{0.5em}

\setlist[itemize]{
    topsep=0.3em,
    itemsep=0.2em
}

\setlist[enumerate]{
    topsep=0.3em,
    itemsep=0.2em
}

\clubpenalty=10000
\widowpenalty=10000

% ------------------------------------------------------------
% ENCABEZADOS Y PIES DE PÁGINA
% ------------------------------------------------------------
\pagestyle{fancy}
\fancyhf{}

\fancyhead[L]{\nouppercase{\leftmark}}
\fancyhead[R]{\materia}

\fancyfoot[C]{\thepage}

\renewcommand{\headrulewidth}{0.4pt}
\renewcommand{\footrulewidth}{0pt}

\fancypagestyle{plain}{
    \fancyhf{}
    \fancyfoot[C]{\thepage}
    \renewcommand{\headrulewidth}{0pt}
}

% ------------------------------------------------------------
% COMANDOS MATEMÁTICOS AUXILIARES
% ------------------------------------------------------------
\newcommand{\abs}[1]{\left|#1\right|}
\newcommand{\norm}[1]{\left\lVert#1\right\rVert}
\newcommand{\vect}[1]{\mathbf{#1}}
\newcommand{\deriv}[2]{\frac{d#1}{d#2}}
\newcommand{\pderiv}[2]{\frac{\partial#1}{\partial#2}}

% ------------------------------------------------------------
% CAJAS ACADÉMICAS SEMÁNTICAS
% ------------------------------------------------------------
\newtcolorbox{definition}[1]{
    enhanced,
    breakable,
    title={Definición: #1},
    fonttitle=\bfseries,
    colback=blue!3,
    colframe=blue!50!black,
    boxrule=0.8pt,
    arc=2mm
}

\newtcolorbox{important}{
    enhanced,
    breakable,
    title={Importante},
    fonttitle=\bfseries,
    colback=yellow!8,
    colframe=orange!70!black,
    boxrule=0.8pt,
    arc=2mm
}

\newtcolorbox{example}[1]{
    enhanced,
    breakable,
    title={Ejemplo: #1},
    fonttitle=\bfseries,
    colback=green!4,
    colframe=green!40!black,
    boxrule=0.8pt,
    arc=2mm
}

\newtcolorbox{formula}[1]{
    enhanced,
    breakable,
    title={Fórmula / Regla: #1},
    fonttitle=\bfseries,
    colback=purple!4,
    colframe=purple!50!black,
    boxrule=0.8pt,
    arc=2mm
}

\newtcolorbox{exercise}[1]{
    enhanced,
    breakable,
    title={Ejercicio: #1},
    fonttitle=\bfseries,
    colback=gray!5,
    colframe=gray!60!black,
    boxrule=0.8pt,
    arc=2mm
}

\newtcolorbox{note}{
    enhanced,
    breakable,
    title={Nota},
    fonttitle=\bfseries,
    colback=white,
    colframe=gray!50,
    boxrule=0.6pt,
    arc=2mm
}

% ============================================================
% DOCUMENTO
% ============================================================
\begin{document}

% ------------------------------------------------------------
% PORTADA
% ------------------------------------------------------------
\begin{titlepage}

\thispagestyle{empty}

\begin{center}

\vspace*{1cm}

{\large \textsc{\universidad}}

\vspace{3cm}

{\Huge \textbf{\materia}}

\vspace{1cm}

{\LARGE \titulo}

\vspace{0.8cm}

\textit{\reflexionmotivacional}

\vspace{1.2cm}

\rule{0.70\textwidth}{0.5pt}

\vspace{0.5cm}

{\large Apuntes de estudio}

\vspace{0.5cm}

\rule{0.70\textwidth}{0.5pt}

\vfill

{\large \autor}

\vspace{0.3cm}

{\large \anio}

\end{center}

\vfill

\begin{flushleft}
\textit{Este es un modesto aporte a los estudiantes de ``\materia'' no pretende sustituir el material oficial de la materia.}
\end{flushleft}

\end{titlepage}

% ------------------------------------------------------------
% ÍNDICE
% ------------------------------------------------------------
\tableofcontents

% ============================================================
% BLOQUE I
% ============================================================
\chapter{Acceso y gestión digital del expediente}

\section{El expediente electrónico y la audiencia virtual}

\subsection{Cómo se ingresa a un expediente}

El acceso a una causa judicial para su consulta o seguimiento se realiza a través del sistema de consulta pública de expedientes. Tratándose de un expediente de naturaleza civil, el ingreso se efectúa seleccionando la Cámara correspondiente —en el caso analizado, la Cámara Nacional Civil— e indicando el año de inicio del expediente. A modo de ejemplo de uso del sistema, se trabajó sobre el expediente identificado con el número 6923/25.

Resulta importante tener presente que una audiencia no constituye un acto aislado ni espontáneo: es la consecuencia de una serie de actos procesales previos —demanda, contestación de demanda, entre otros— que desembocan en la fijación de la fecha de audiencia por parte del juzgado.

\begin{note}
La notificación de la fecha de audiencia se realiza mediante una \textbf{cédula de tipo electrónica}, dirigida al domicilio procesal constituido por cada parte. En la práctica del fuero, la audiencia preliminar es identificada coloquialmente como \textbf{«audiencia de 360»}.
\end{note}

\subsection{El acceso a las salas virtuales}

Las audiencias virtuales son, en principio, de carácter público. Sin embargo, el ingreso efectivo a la sala virtual requiere contar con un enlace de acceso que otorga el propio juzgado; sin ese enlace, no es posible ingresar por cuenta propia, aun tratándose de un acto formalmente público.

\subsection{La importancia de la calidad de imagen y sonido en la audiencia virtual}

La modalidad virtual de las audiencias exige un estándar de calidad técnica que no es menor: dado que una audiencia puede ser eventualmente revisada por la instancia de Cámara, resulta necesario que la imagen y el sonido permitan una correcta visualización de cada interviniente. Lo esperable es que cada parte se encuentre en un entorno adecuado —por ejemplo, sentada frente a un escritorio— de modo tal que pueda ser vista con claridad durante todo el acto.

\subsection{La solemnidad del acto judicial y su raíz histórica}

El lugar y la forma en que se desarrolla un acto judicial no son un detalle accesorio. Los espacios cumplen una función simbólica además de práctica: así como un templo predispone a una actitud determinada frente a una plaza, el ámbito de los tribunales —con sus sellos, sus formas y su disposición— genera un marco de significado propio. Cuando ese marco formal se pierde y el acto se vuelve excesivamente doméstico, se debilita también la solemnidad que históricamente ha caracterizado a la administración de justicia.

\begin{example}{La traditio romana como antecedente de la solemnidad jurídica}
En el derecho romano, ante la ausencia generalizada de alfabetización y de instrumentos adecuados para dejar constancia escrita de los actos jurídicos, se recurría a ceremonias simbólicas para asegurar que un acto trascendente fuera recordado con el paso del tiempo. En la compraventa de un terreno, por ejemplo, el vendedor entregaba al comprador un pedacito de tierra —la \textit{traditio}— y luego se la quitaba de manera simbólica, ante la presencia de dos o tres testigos. La función de esa ceremonia era generar un recuerdo duradero del acto en quienes presenciaban la escena, de un modo comparable a como se recuerda cualquier acto trascendente en la vida de una persona.
\end{example}

La conclusión pedagógica que conecta la audiencia virtual contemporánea con la ceremonia romana es la siguiente: cuando el acto judicial pierde su marco formal, se debilita también su capacidad de ser recordado, comprendido y respetado por quienes participan de él.

% ============================================================
% BLOQUE II
% ============================================================
\chapter{La prescripción adquisitiva y la economía del proceso}

\section{Concepto, presupuestos y alternativas de la usucapión}

\subsection{Concepto de usucapión}

\begin{definition}{Usucapión (prescripción adquisitiva)}
La usucapión es un derecho real por el cual el transcurso del tiempo permite que una persona que ha poseído una cosa como dueño solicite a un juez el reconocimiento de esa titularidad, aun cuando no haya sido su titular registral originario, en virtud de haber ejercido los derechos correspondientes —por ejemplo, pagando los impuestos y comportándose como propietario— durante el plazo que la ley exige.
\end{definition}

\begin{formula}{Código Civil y Comercial de la Nación — Artículo 1897}
«La prescripción para adquirir es el modo por el cual el poseedor de una cosa adquiere un derecho real sobre ella, mediante la posesión durante el tiempo fijado por la ley.»

\textit{CCyC, Libro Cuarto, Título I, Capítulo 2.}
\end{formula}

\begin{formula}{Código Civil y Comercial de la Nación — Artículo 1899}
«Si no existe justo título o buena fe, el plazo es de veinte años [\ldots].»

\textit{CCyC, Libro Cuarto, Título I, Capítulo 2 (prescripción larga o veinteañal).}
\end{formula}

\subsection{La contrapartida: qué puede alegar la parte demandada}

Frente a una demanda de usucapión, la parte demandada puede oponer defensas tales como que no se cumplió el plazo legal, que la posesión invocada no existió tal como se relata, o que dicha posesión fue interrumpida en algún momento.

Lo que persigue la parte actora, en definitiva, es que el juez le reconozca la titularidad dominial y la convierta en propietaria, como si hubiera adquirido el bien mediante escritura pública. De resultar favorable la sentencia, esa titularidad queda inscripta en el Registro de la Propiedad Inmueble a nombre del actor, quien podrá disponer del bien como si lo hubiese comprado.

\begin{note}
La usucapión es una figura de uso frecuente en la Argentina, particularmente en el interior del país —en campos, chacras y pequeños terrenos—. En la Ciudad de Buenos Aires es relativamente menos habitual, debido a un mayor grado de organización registral.
\end{note}

\subsection{Un caso ilustrativo: la mediación imposible}

\begin{example}{El campo cercano a La Plata: cuando la mediación no alcanza}
Un campo próximo a La Plata era ocupado desde hacía aproximadamente cuarenta años por personas que desarrollaban allí un emprendimiento de cultivo y que habían pagado casi todos los impuestos correlativos —uno de los elementos que más acredita la usucapión—. Dado el tiempo transcurrido, resultaba prácticamente imposible que los titulares registrales originales continuaran con vida. Ante la pregunta de si, en un caso así, correspondería una mediación previa, la respuesta fue negativa: la usucapión no es un derecho disponible por acuerdo de partes. Aun si se pudiera reunir a eventuales herederos o titulares para negociar, cualquier transferencia de la titularidad por acuerdo exigiría otro acto jurídico distinto —una donación o una venta—, que no reemplaza la necesidad de una declaración judicial que constate el cumplimiento de los requisitos legales de la usucapión.
\end{example}

La conclusión conceptual de este caso es que la usucapión no es un derecho negociable entre partes: requiere necesariamente una \textbf{declaración judicial}, porque nace del cumplimiento objetivo de requisitos legales —posesión y tiempo— y no de la voluntad de quien figura como titular registral, muchas veces inexistente o inhallable.

\subsection{Usucapión versus juicio de escrituración}

Cuando el actor cuenta con un boleto de compraventa, puede optar entre promover una acción de usucapión o un juicio de escrituración; ambas vías pueden conducir al mismo resultado práctico —la obtención del título de propiedad—. La elección entre una u otra depende de una cuestión estratégica vinculada a la documentación y a la prueba disponible en cada caso concreto.

\begin{table}[H]
\centering
\caption{Usucapión y juicio de escrituración: criterios comparativos}
\begin{tabularx}{\textwidth}{X X}
\toprule
\textbf{Usucapión (prescripción adquisitiva)} & \textbf{Juicio de escrituración} \\
\midrule
Se funda en la posesión y el paso del tiempo (20 años), con independencia de que exista o no un boleto de compraventa. &
Se funda en el boleto de compraventa; exige acreditar la obligación de escriturar asumida por el vendedor. \\
\addlinespace
Resultado depende de la prueba de los actos posesorios (pago de impuestos, expensas, uso, testigos). &
Resultado depende de la prueba del vínculo contractual y del incumplimiento de la obligación de escriturar. \\
\addlinespace
Resulta especialmente útil cuando el vendedor original desapareció, se disolvió o es inhallable. &
Requiere, en principio, poder dirigir la acción contra quien se obligó a escriturar (o sus sucesores). \\
\addlinespace
Es la vía elegida cuando se considera contar con todos los elementos para probar la posesión. &
Es la vía elegida cuando la relación contractual y la identidad del obligado están más claras. \\
\bottomrule
\end{tabularx}
\end{table}

\section{Tasa de justicia, bono y caja previsional}

\subsection{La tasa de justicia y el bono de derecho fijo}

Todo proceso judicial exige el pago de la tasa de justicia y del bono de derecho fijo, cuya forma de abono varía según la jurisdicción.

En la Ciudad Autónoma de Buenos Aires, el bono se abona a través de una plataforma del Banco Ciudad de Buenos Aires: se ingresa a la plataforma, se selecciona la opción de pago del bono, se indica la Cámara o fuero correspondiente, el número de expediente y el CUIT del letrado, y el pago se efectúa de manera virtual mediante transferencia, acreditándose directamente en el expediente sin necesidad de acompañar comprobante adicional.

\begin{note}
En la provincia de Buenos Aires, el pago se continúa gestionando a través del Colegio de Abogados (CRA): el comprobante se descarga en formato PDF y se acompaña junto con el resto de la documentación de la demanda.
\end{note}

\subsection{Las cajas de previsión profesional}

En la provincia de Buenos Aires, y en la mayoría de las provincias argentinas, existe una caja de previsión social específica para abogados, independiente de la caja pública general. Estas cajas son de afiliación obligatoria: un porcentaje de los honorarios profesionales —entre el 10\% y el 15\%— se destina a dicha caja, con aportes a cargo tanto del propio profesional como, en parte, de la contraparte. El incumplimiento de este aporte impide obtener regulaciones judiciales de honorarios o la expedición de testimonios, entre otras consecuencias.

\begin{important}
Esta obligación no es uniforme en todo el país. En la Capital Federal, los abogados pueden inscribirse como autónomos ante la AFIP y no existe una caja previsional especial que acreditar, más allá del pago del bono. En cambio, en la provincia de Buenos Aires, en Tierra del Fuego y en otras provincias —como Córdoba— que cuentan con caja de profesionales independiente, el aporte previsional es obligatorio. Si al cierre del año no se alcanzó, mediante los depósitos de honorarios, el monto mínimo exigido, el profesional debe aportar la diferencia igualmente para mantener vigente tanto la matrícula como la afiliación a la caja.
\end{important}

Esta distinción resulta clave para cualquier estudiante que planee ejercer la profesión: el costo económico de litigar, y las obligaciones profesionales asociadas, varían según la jurisdicción en la que se actúe, y deben calcularse antes de iniciar cualquier proceso.

% ============================================================
% BLOQUE III
% ============================================================
\chapter{Oralidad y audiencia preliminar}

\section{El protocolo de oralidad en el proceso civil}

\subsection{Origen del protocolo de oralidad}

\begin{definition}{Protocolo de oralidad}
Es una política impulsada desde el Ministerio de Justicia con anterioridad a la pandemia, orientada a aplicar en el proceso civil, en la mayor medida posible, la producción oral de la prueba —especialmente la prueba pericial— durante la audiencia preliminar, en lugar de resolver la totalidad de las cuestiones probatorias por escrito.
\end{definition}

\subsection{Por qué la prueba oral resulta más clara que la escrita}

Entre las distintas pruebas que pueden producirse oralmente, la pericial ocupa un lugar especialmente valioso: no es lo mismo que el perito presente su informe únicamente por escrito, a que el juez pueda preguntarle directamente, en la audiencia, cómo se produjo determinado hecho técnico.

\begin{example}{El chalet con la filtración: la claridad de la palabra oral}
En un juicio relativo a un problema de filtración en un chalet, no resultaba claro si el defecto obedecía a un error de construcción o a un problema con las tejas empleadas. La demanda había sido dirigida contra todos los posibles responsables, precisamente por esa incertidumbre. En la audiencia, el juez interrogó directamente al perito, quien explicó oralmente que el problema se había originado en el momento de la colocación de las tejas, y no en su fabricación, lo que permitió identificar con rapidez a quién correspondía la responsabilidad. Esa misma explicación, presentada únicamente por escrito, habría resultado mucho más extensa, discutible y difícil de comprender. Tras la explicación oral del perito, las partes dejaron de tener motivos para seguir discutiendo, porque el resultado del proceso ya resultaba previsible.
\end{example}

\subsection{Ventajas y exigencias del protocolo de oralidad}

De aplicarse plenamente, el protocolo de oralidad permitiría reducir drásticamente —a una fracción estimada del 20\% o menos— el tiempo que insumen los procesos tramitados íntegramente por escrito, además de acercar al juez a los hechos reales y limitar las maniobras dilatorias.

\begin{important}
La oralidad tiene una contracara: exige que las partes y sus letrados lleguen a la audiencia preparados para responder preguntas en vivo, sin poder ampararse en la elaboración diferida de un escrito. Quien no tiene claro qué sostiene, qué pretende y qué prueba ofrece, queda expuesto frente al tribunal.
\end{important}

\subsection{La audiencia de vista de causa en las provincias}

La denominada «audiencia de 360» —tal como se la conoce coloquialmente en el fuero nacional— guarda equivalencia con la «audiencia de vista de causa» prevista en los códigos procesales de varias provincias, aun cuando no lleve formalmente ese nombre en la Nación. El protocolo de oralidad, diseñado originalmente para el ámbito nacional, resulta en lo esencial similar al que rige en distintas jurisdicciones provinciales.

\section{Audiencia preliminar en vivo: caso de la cochera}

\subsection{Apertura de la audiencia}

La audiencia preliminar observada correspondió al expediente caratulado \textit{«Chester, Claudio c/ Construcciones Mirage S.A. s/ prescripción adquisitiva»}. Al momento de la apertura, se dejó constancia de que no había comparecido representante alguno del Gobierno de la Ciudad —citado por el interés fiscal comprometido— ni de la parte demandada.

\subsection{El interrogatorio a la parte actora}

El actor relató que en 1989 había adquirido, mediante boleto de compraventa, la cochera identificada como unidad funcional 28 del edificio sito en Independencia 1452.
% TODO: contenido incompleto o ambiguo en las notas originales (el piso indicado en la transcripción es "64", dato que no resulta verosímil y que se conserva tal como figura en el apunte).
Pagó la totalidad del precio, pero nunca logró que se le escriturara la unidad; desde entonces la usó personalmente durante un tiempo, y luego la alquiló bajo contratos siempre verbales, abonando en todo momento los gastos correspondientes a la cochera.

Consultado sobre la constructora demandada, el actor indicó no tener certeza sobre si otras unidades del edificio atravesaban una situación similar a la suya. Sobre su relación con el consorcio, señaló que no participaba de las asambleas —por tratarse de una cochera de escaso interés para él en ese aspecto— pero que siempre se había sentido, a los efectos prácticos, como propietario del bien. Relató además haber insistido en su momento para lograr la escrituración, sin éxito, y aclaró que no tiene intención de vender la unidad, sino de regularizar su situación dominial.

\subsection{Los testigos ofrecidos y su vinculación con los hechos}

Antes de admitir formalmente a los testigos propuestos por la parte actora —Alberto Damián Cabuli y María Teresa Gamboa—, la jueza examinó su utilidad concreta para la causa.

\begin{important}
El criterio de admisión de la prueba testimonial no se limita a la mera oferta del testigo: el objeto de la audiencia preliminar es que la prueba propuesta esté efectivamente dirigida a acreditar los hechos controvertidos —en este caso, la posesión sostenida durante más de veinte años—. No corresponde citar a una persona que no pueda aportar nada relevante al proceso.
\end{important}

El actor explicó que uno de los testigos lo había acompañado en trámites vinculados a la cochera y, en una oportunidad, se la había alquilado; la segunda testigo trabajaba junto al primero. Sobre esa base, ambos testimonios fueron admitidos por su vinculación con los hechos a probar.

\subsection{Producción de la prueba y prueba desistida}

La parte actora mantuvo la prueba testimonial y la informativa —dirigida al Gobierno de la Ciudad, a Aguas Argentinas y al consorcio de propietarios—, pero desistió de la prueba pericial caligráfica y de la inspección ocular, por considerarlas redundantes: el desinterés ya manifestado por el Gobierno de la Ciudad y la prueba informativa ordenada cumplían, a criterio de la parte, la misma función.

Se dejó además constancia de un dato relevante para el desarrollo posterior del proceso: la firmante del boleto de compraventa en representación de Construcciones Mirage S.A. actuaba como representante de una sociedad que ya se encontraba disuelta desde el año 2001.

\begin{formula}{Acta de la audiencia preliminar (síntesis)}
«Buenos Aires, a los 18 días de septiembre de 2026 [\ldots] se llama a la audiencia convocada en este proceso y comparece por la parte actora el señor Claudio Chester, asistido por su letrada [\ldots]. Se resuelve abrir la causa a prueba. Se ordena prueba informativa a Aguas Argentinas, al Gobierno de la Ciudad y al consorcio propietario [\ldots]. La testimonial que se ordena será llevada a cabo el día de la audiencia de vista de causa [\ldots], quienes deberán comparecer de manera presencial en la sede del juzgado.»

\textit{Transcripción de lo dictado por la jueza durante la audiencia, reconstruida a partir del audio de la clase.}
\end{formula}

\section{Análisis posterior: prueba y estrategia procesal}

\subsection{Primeras impresiones de la clase}

El estilo de conducción de la audiencia por parte de la jueza fue valorado como poco frecuente: mantuvo en todo momento el control del acto, formulando directamente al actor todas las preguntas necesarias para reconstruir los hechos. Se señaló que, de tratarse de un caso más complejo que el analizado —donde el actor contaba con boleto, pago de expensas y un único punto débil, la falta de escrituración—, la magistrada probablemente habría adoptado una postura interrogativa aún más incisiva.

\subsection{Un error común: ofrecer testigos sin explicar para qué}

\begin{important}
Constituye un error frecuente en la práctica profesional ofrecer un testigo sin precisar, ante el juez, para qué hecho concreto se lo propone. No basta con indicar que una persona será citada como testigo: es necesario explicar qué hecho específico puede acreditar cada uno de los testigos ofrecidos, de manera que el tribunal no quede sin saber cuál es la utilidad probatoria de cada testimonio.
\end{important}

\subsection{Por qué se descartó parte de la prueba ofrecida}

Cuando una prueba puede ser suplida por otra ya ordenada, resulta razonable desistir de ella para no incurrir en un gasto de tiempo innecesario. En el caso analizado, se desistió de la prueba pericial y de la inspección ocular, en la medida en que la prueba informativa ya dispuesta cumplía una función equivalente. Se mencionó, además, que en ciertos casos la inspección ocular puede sustituirse por medios audiovisuales, en la medida en que ello resulte útil y comprensible para el proceso.

\subsection{El desafío que queda planteado: una demandada disuelta}

\begin{example}{Una sociedad que ya no existe}
La parte demandada, Construcciones Mirage S.A., se encontraba disuelta desde el año 2001, y el Gobierno de la Ciudad —citado por el interés fiscal comprometido— había manifestado expresamente su desinterés en el litigio. Ante esta situación, se planteó el interrogante de cómo continuar el proceso frente a una contraparte que, en los hechos, ya no existe, y qué gestiones corresponde realizar al letrado de la parte actora para sortear ese obstáculo procesal.
\end{example}

Este interrogante —cómo obtener información institucional cuando no hay una contraparte activa a quien reclamarle— da paso al desarrollo de la figura del oficio judicial.

% ============================================================
% BLOQUE IV
% ============================================================
\chapter{Los oficios judiciales}

\section{El artículo 400 del CPCCN y su diligenciamiento electrónico}

\subsection{Qué es un oficio y a quién se dirige}

\begin{definition}{Oficio judicial}
El oficio es un pedido de comunicación o de información que se dirige siempre a una entidad —pública o privada— y nunca a una persona física, con el fin de que los datos obrantes en los registros de esa institución sean puestos a disposición del proceso.
\end{definition}

\begin{note}
La distinción entre oficio y prueba testimonial es clara: cuando se busca obtener información de una \textbf{entidad}, corresponde dirigirle un oficio; cuando la información que se pretende acreditar depende del conocimiento de una \textbf{persona física}, corresponde ofrecerla —o a quien la conozca— como testigo.
\end{note}

\subsection{El oficio del artículo 400 y el oficio 3003}

Existen dos grandes categorías de oficios según su finalidad. El primero, denominado \textbf{oficio 3003}, se emplea de manera específica en procesos sucesorios o concursales: antes de promover un proceso sucesorio, una quiebra o cualquier proceso universal, se dirige este oficio al Registro de Procesos Universales, a fin de verificar si ya existe iniciada una sucesión o una quiebra a nombre de la persona o empresa de que se trate.

El segundo tipo, de uso general, es el regulado por el artículo 400 del Código Procesal Civil y Comercial de la Nación.

\begin{formula}{Código Procesal Civil y Comercial de la Nación — Artículo 400 (Atribuciones de los letrados patrocinantes)}
«Los pedidos de informes, testimonios y certificados, así como los de remisión de expedientes ordenados en el juicio, serán requeridos por medio de oficios firmados, sellados y diligenciados por el letrado patrocinante, con transcripción de la resolución que los ordena y que fija el plazo en que deberán remitirse.»

\textit{CPCCN, Libro Segundo, Título II, Capítulo V, Sección 3.ª}
\end{formula}

\begin{note}
Con la implementación del expediente digital, buena parte de las exigencias formales del artículo 400 —como la firma y el sellado— se ha simplificado considerablemente, y solo cobra relevancia práctica cuando la entidad destinataria no figura entre las disponibles en el sistema de gestión judicial.
\end{note}

\subsection{Diligenciamiento electrónico del oficio}

El módulo de oficios del sistema de gestión judicial requiere completar, en orden, los siguientes campos: la jurisdicción y el número de expediente; la entidad destinataria (por ejemplo, una entidad bancaria); el tipo de oficio (artículo 400); el texto del despacho judicial que ordena la medida; y la oficina de recepción del destinatario. Una vez completado el trámite, el oficio queda cargado directamente en el expediente como diligenciado, sin necesidad de presentar un escrito adicional solicitando al juzgado que lo diligencie, tal como se exigía bajo el procedimiento anterior en soporte papel.

\subsection{Cuándo el destinatario no está en el sistema: el oficio en papel}

Cuando la entidad destinataria no figura en la lista de organismos disponibles en el sistema —por ejemplo, determinadas empresas de comercio electrónico—, corresponde volver al procedimiento tradicional en soporte papel.

\begin{formula}{Estructura del oficio en papel (dictado en clase)}
«Tengo el agrado de dirigirme a usted en mi calidad de letrado patrocinante [o apoderado] de la parte actora [o demandada] en los autos caratulados «\ldots», que tramitan por ante el Juzgado [\ldots], a cargo de [\ldots], Secretaría a cargo de [\ldots], con domicilio en [\ldots], a fin de solicitarle quiera tener a bien informar [el objeto del pedido] [\ldots]. [Se transcribe a continuación el auto que ordena la medida]. Lo saludo atentamente.»

\textit{Modelo dictado durante la clase, a completar con los datos de cada expediente.}
\end{formula}

\begin{important}
Al transcribir la resolución judicial dentro del cuerpo del oficio, no corresponde copiar el auto en su integridad, sino únicamente el fragmento operativo en el que se dispone la medida —por ejemplo, la parte que ordena «líbrese el oficio solicitado»— o la fórmula equivalente que utilice la resolución.
\end{important}

Cuando la respuesta al oficio en papel se recibe también en soporte papel —sin recibo digital que acredite el trámite— el procedimiento para acreditar el diligenciamiento consiste en confeccionar dos copias del oficio: una que queda en poder del letrado y otra que se presenta ante la institución destinataria para obtener el sello de recibido. Con esa copia sellada, se digitaliza el documento y se presenta un escrito ante el expediente electrónico manifestando que se acredita el diligenciamiento del oficio, acompañando la copia digitalizada correspondiente.

\subsection{Presentación del escrito y cierre de la clase}

La presentación de este escrito ante el expediente electrónico exige seleccionar el expediente principal —y no un eventual incidente con la misma carátula—, elegir el tipo de escrito correspondiente (habitualmente, «mero trámite») y adjuntar en un único archivo PDF tanto el escrito como la copia del oficio diligenciado. Completada esta carga, el escrito queda registrado en el expediente como enviado.

\begin{note}
Sobre la modalidad de evaluación de la materia, se adelantó que consistirá en tres o cuatro preguntas a desarrollar de manera simple, pudiendo incluir consignas para redactar un escrito sencillo o para narrar y ofrecer como prueba un conjunto de hechos determinados.
\end{note}

% ============================================================
% CUADRO CORNELL Y SÍNTESIS FINAL
% ============================================================
\chapter{Cuadro Cornell y síntesis final}

El siguiente cuadro reúne, según el método Cornell, las preguntas centrales que permiten repasar activamente los contenidos desarrollados en esta clase.

\begin{longtable}{
    >{\raggedright\arraybackslash}p{0.30\textwidth}
    >{\raggedright\arraybackslash}p{0.60\textwidth}
}
\toprule
\textbf{Preguntas / palabras clave} & \textbf{Notas / respuestas} \\
\midrule
\endfirsthead

\toprule
\textbf{Preguntas / palabras clave} & \textbf{Notas / respuestas} \\
\midrule
\endhead

\midrule
\multicolumn{2}{r}{\textit{Continúa en la página siguiente}} \\
\endfoot

\bottomrule
\endlastfoot

\textbf{¿Cómo se accede a un expediente electrónico?} &
A través del sistema de consulta pública, ingresando por «expediente», seleccionando la Cámara (por ejemplo, la Cámara Nacional Civil) y el año del expediente. Las audiencias virtuales son en principio públicas, pero requieren el enlace de acceso que otorga el juzgado. \\
\addlinespace

\textbf{¿Por qué importa la solemnidad del lugar desde donde se litiga?} &
Porque el marco formal del acto judicial ayuda a que las partes lo recuerden, lo comprendan y lo respeten. Se lo relacionó con la \textit{traditio} romana: una ceremonia simbólica con testigos que aseguraba que un acto trascendente fuera recordado con el paso del tiempo. \\
\addlinespace

\textbf{¿Qué es la prescripción adquisitiva (usucapión)?} &
Es el modo por el cual el poseedor de una cosa adquiere un derecho real sobre ella mediante la posesión durante el tiempo fijado por la ley (art.\ 1897 CCyC). Sin justo título ni buena fe, el plazo es de veinte años (art.\ 1899 CCyC). Si la sentencia es favorable, el derecho se inscribe en el Registro de la Propiedad Inmueble. \\
\addlinespace

\textbf{¿Por qué la usucapión no se puede resolver por mediación?} &
Porque no es un derecho disponible por acuerdo de partes: requiere una declaración judicial que verifique el cumplimiento de los requisitos legales (posesión y tiempo). Transferir la titularidad por acuerdo exigiría otro acto jurídico, como una donación o una venta. \\
\addlinespace

\textbf{Usucapión vs.\ juicio de escrituración: ¿cuándo conviene cada uno?} &
Ambos pueden llevar al mismo resultado (obtener el título de propiedad). La elección es estratégica: depende de si el actor cuenta con mejor prueba de la posesión (usucapión) o de la relación contractual y el incumplimiento de escriturar (juicio de escrituración). \\
\addlinespace

\textbf{¿Qué hay que pagar para litigar y quién lo cobra?} &
La tasa de justicia y el bono de derecho fijo (en CABA, pagado digitalmente a través del Banco Ciudad; en Provincia, por el Colegio de Abogados). Además, en la mayoría de las provincias existe una caja de previsión profesional obligatoria (10-15\% de los honorarios); en la Capital Federal, los abogados se inscriben como autónomos ante la AFIP. \\
\addlinespace

\textbf{¿Qué es el protocolo de oralidad?} &
Una política, impulsada desde el Ministerio de Justicia antes de la pandemia, para aplicar en el proceso civil, en la mayor medida posible, la producción oral de la prueba (especialmente la pericial) durante la audiencia preliminar, en lugar de resolver todo por escrito. \\
\addlinespace

\textbf{¿Qué ventajas y qué exigencias trae la oralidad?} &
Reduce drásticamente el tiempo del proceso, acerca al juez a los hechos reales y limita las maniobras dilatorias. A cambio, exige que las partes y sus letrados lleguen preparados para responder preguntas en vivo, sin poder ocultarse detrás de un escrito. \\
\addlinespace

\textbf{En la audiencia observada, ¿qué debía probar el actor?} &
Los actos posesorios sostenidos durante veinte años sobre una cochera adquirida por boleto de compraventa en 1989: uso, pago de expensas, comportamiento como dueño, y testigos vinculados efectivamente a esos hechos (no meros conocidos). \\
\addlinespace

\textbf{¿Qué prueba se desistió en la audiencia y por qué?} &
La pericial caligráfica y la inspección ocular, porque la prueba informativa ya ordenada (y la manifestación de desinterés del Gobierno de la Ciudad) cumplía la misma función: ofrecer prueba redundante implica un gasto de tiempo innecesario. \\
\addlinespace

\textbf{¿Qué error señaló el profesor sobre el ofrecimiento de testigos?} &
No basta con ofrecer un testigo: hay que explicar para qué se lo ofrece y qué hecho concreto puede acreditar, porque el juez solo admitirá testigos con vinculación real y útil a los hechos que se buscan probar. \\
\addlinespace

\textbf{¿Qué es un oficio y en qué se diferencia de un testigo?} &
Es un pedido de información dirigido a una entidad pública o privada (nunca a una persona física), regulado por el artículo 400 del CPCCN. Para obtener información sobre una persona, en cambio, corresponde ofrecerla —o a quien la conozca— como testigo. \\
\addlinespace

\textbf{¿Qué es el oficio 3003?} &
Un oficio específico, previo al inicio de un proceso sucesorio o concursal, dirigido al Registro de Procesos Universales, para verificar si ya existe una sucesión o una quiebra iniciada respecto de determinada persona o empresa. \\
\addlinespace

\textbf{¿Cómo se diligencia hoy un oficio del artículo 400?} &
De manera electrónica: se selecciona el expediente, la entidad destinataria, el tipo de oficio, el texto del despacho judicial y la oficina de recepción; el sistema lo carga automáticamente al expediente, sin necesidad de presentarlo en papel, salvo que el destinatario no figure en el sistema. \\
\addlinespace
\midrule

\multicolumn{2}{p{0.90\textwidth}}{
\textbf{Resumen:} La clase reconstruyó, de punta a punta, el ciclo completo de un litigio civil concreto. Comenzó por lo más operativo —cómo ingresar a un expediente electrónico y a una audiencia virtual— y de ahí derivó hacia la pregunta de fondo del caso observado: qué es la prescripción adquisitiva, qué debe probarse para obtenerla y por qué exige siempre una sentencia judicial, a diferencia de una negociación privada. En el camino aparecieron los costos concretos de litigar —tasa de justicia, bono, caja previsional— que todo abogado debe calcular antes de iniciar una demanda. El núcleo de la clase fue la audiencia preliminar en vivo, entendida como la expresión práctica del protocolo de oralidad: un espacio donde la jueza interrogó directamente al actor para reconstruir veinte años de posesión sobre una cochera, admitió o descartó prueba según su utilidad real, y dejó planteado un problema procesal —una sociedad demandada ya disuelta— que la clase resolvió recurriendo a la figura del oficio judicial (art.\ 400 CPCCN), diferenciándolo con precisión de la prueba testimonial y describiendo su diligenciamiento tanto electrónico como, subsidiariamente, en papel. El hilo conductor de toda la clase fue, en definitiva, un mismo principio: la eficacia de un proceso judicial depende tanto de conocer la norma como de saber para qué sirve cada prueba, cada trámite y cada gesto formal dentro de él.
} \\

\end{longtable}

\section*{Normativa citada en estos apuntes}
\addcontentsline{toc}{section}{Normativa citada en estos apuntes}

\begin{itemize}
    \item Código Civil y Comercial de la Nación, artículo 1897 (concepto de prescripción adquisitiva).
    \item Código Civil y Comercial de la Nación, artículo 1899 (plazo de veinte años sin justo título ni buena fe).
    \item Código Procesal Civil y Comercial de la Nación, artículo 400 (atribuciones de los letrados patrocinantes respecto de los oficios).
\end{itemize}

\end{document}
